% !TeX root = ../../tesis.tex

% =============================================================================
%  4.2.3 — ENTRADAS Y VALIDACIÓN
% =============================================================================
\section{Entradas y Validación de Datos}
\label{sec:vft-entradas}

\subsection{Formato de Ingesta: GeoJSON}
\label{subsec:vft-geojson}

VFTModel se alimenta de los datos geográficos producidos por Apimetro en formato
\geojson{} \citep{butler2016geojson}. Cada petición devuelve una
\texttt{FeatureCollection} con dos tipos de elementos:

\begin{itemize}
\item \textbf{Estaciones} (campo \texttt{tipo\_entidad} = \texttt{estacion}):
representadas
    como puntos geográficos (\texttt{Point}), con coordenadas GPS y atributos
    operativos —nombre, sistema, jerarquía de transporte, alcaldía y si es un
    nodo \textsc{cetram}—.
\item \textbf{Rutas} (campo \texttt{tipo\_entidad}\allowbreak{} = \texttt{ruta}
): representadas como
    trazos lineales (\texttt{MultiLineString}), con las métricas que el motor de
    impedancia requiere —velocidad operacional, frecuencia de servicio, tipo de
    derecho de vía y sentido de circulación—.
\end{itemize}

Antes de procesar cualquier elemento, VFTModel valida que todos los campos
requeridos existan y sean válidos, como se describe en la siguiente sección.

\subsection{Validación Estricta con Pydantic}
\label{subsec:vft-pydantic}

Cada \texttt{Feature} recibido es sometido a una validación estricta mediante
\textbf{Pydantic}, una biblioteca de Python que verifica que los datos cumplan
exactamente los contratos de tipo definidos en el modelo. Si algún campo es
inválido —por ejemplo, un sistema de transporte no reconocido o un valor de
velocidad negativo—, el sistema lo rechaza antes de que contamine los cálculos.
Este filtro está implementado en \texttt{src/api/schemas/schemas.py}.

El contrato verifica ocho campos críticos: \texttt{sistema} (Enum de 12
valores), \texttt{tipo\_entidad} (\texttt{estacion} o \texttt{ruta}),
\texttt{jerarquia\_transporte}, \texttt{derecho\_de\_via},
\texttt{velocidad\_promedio\_kmh}, \texttt{frecuencia\_minutos},
\texttt{sentido} y \texttt{es\_cetram}. La especificación completa del contrato
de datos se presenta en el Cuadro~\ref{tab:vft-campos-validados} del Anexo~
\ref{anx:vft-validacion}.

\subsection{Catálogo de Sistemas de Transporte}
\label{subsec:vft-sistemas}

VFTModel reconoce exactamente doce sistemas de transporte metropolitano. Cuando
los datos de velocidad o frecuencia no llegan desde Apimetro, el modelo aplica
valores de \termino{fallback} calibrados a la operación real de cada sistema.
Los sistemas estructurales de mayor peso en el análisis son el STC Metro (36
km/h, vía exclusiva), Metrobús (16.3 km/h, vía confinada), Tren Ligero (22 km/h,
vía exclusiva), \textsc{rtp} (14 km/h, vía mixta) y Corredores Concesionados (11
km/h, vía mixta). El catálogo completo con claves de código y parámetros de
\termino{fallback} se presenta en el Cuadro~\ref{tab:vft-sistemas} del Anexo~
\ref{anx:vft-catalogo}.

El tipo de derecho de vía es el parámetro que determina el coeficiente de
fricción aplicado a cada arista del grafo, como se describe en la Sección~
\ref{subsec:vft-cf}.
